\chapter{暴脹、結構與初始條件的消失}

\section{{\proseunicodefont Ω}與平坦性的證據}

一位學生問到大\(\Omega\)，討論因此回到弗里德曼–羅伯森–沃克方程。就現在這個遠離輻射主導時期的宇宙而言，膨脹定律寫成真空能項與物質項；輻射項則因為已經非常小而略去：
\begin{equation}
H^2=\left(\frac{\dot a}{a}\right)^2=\frac{8\pi G}{3}\left(\rho_0+\rho_m\right)-\frac{k}{a^2}.
\end{equation}
真空項不會稀釋，物質項則會：
\begin{equation}
\rho_m=\frac{M}{a^3}.
\end{equation}
\(\Omega\)表示膨脹方程中各種組分所占的比例。

將方程兩邊除以\(H^2\)，各項經過歸一化後的總和便為一。採用經整理的記號，
\begin{equation}
\begin{aligned}
\Omega_0&=\frac{8\pi G\rho_0}{3H^2},\\
\Omega_m&=\frac{8\pi G\rho_m}{3H^2},\\
\Omega_k&=-\frac{k}{a^2H^2},
\end{aligned}
\end{equation}
因此
\begin{equation}
\Omega_0+\Omega_m+\Omega_k=1.
\end{equation}
\editorialnote{這一段的音訊與黑板內容都不太清楚，講者也在幾個名稱之間切換，指稱真空貢獻。此處將方程除以\(H^2\)後，整理了論證所用的記號。}
若空間是平坦的，\(k=0\)，因此
\begin{equation}
\Omega_0+\Omega_m=1.
\end{equation}
物質項與真空項之和若大於一，曲率為正；若小於一，曲率為負。

觀測上的推論很直接：測量普通物質與暗物質，測量哈勃膨脹，並在極大尺度上以幾何方法測量曲率。也要留意，宇宙正在加速。判斷真空項的符號，不能只看\(H\)本身，而要看\(\dot a/a\)趨近常數的趨勢；這表示尺度因子會趨近指數膨脹的形式，
\begin{equation}
a(t)\propto e^{Ht}.
\end{equation}
這些測量彼此交叉檢驗，顯示宇宙非常接近平坦，偏差只有百分之幾。

\begin{classroomqa}
\lecturetimestamp{00:10:50}
\audiencequestion{物質項包括暗物質嗎？}
\lecturerresponse{是。就目前的討論而言，物質項包括暗物質。也可以把輻射加進去，不過如今輻射所占的比例非常小，可以忽略。}
\end{classroomqa}

\begin{classroomqa}
\lecturetimestamp{00:11:11}
\audiencequestion{除了哈勃常數之外，還有其他方法可以衡量暗能量嗎？}
\lecturerresponse{重點不在\(H\)本身的數值，而在於膨脹如何趨近指數定律。再結合物質的測量結果，以及以幾何方法測得的平坦性，就能推論出真空項。}
\end{classroomqa}

\section{從平坦性到結構}

宇宙為什麼如此平坦？一種解釋是暴脹。平坦意味著宇宙在數十億、甚至數百億光年的尺度上非常大、非常平滑，也非常均勻。指數膨脹可以把一個區域拉伸到呈現這種樣貌。

不過，好事也可能做過頭。完全平滑的宇宙會一直保持完全平滑。引力無法從精確的對稱性中產生成團結構。既然現在存在結構，這些結構必定是從早期的微小不均勻性開始形成。將現在的宇宙往回追溯，到了透明時期，不均勻性的大小約為\(10^5\)分之一：
\begin{equation}
\frac{\delta\rho}{\rho}\sim 10^{-5}.
\end{equation}
當時的宇宙極其平滑，但並非完全平滑。

從那時起，引力便開始做出一般房間裡的空氣所受的普通力做不到的事。微小的過密區會把物質吸引過來，附近區域因此變得稀疏；於是欠密區愈加欠密，過密區愈加過密。引力會增強這種反差，而不是將它抹平。

自然會問，為什麼這種情況沒有更早發生。答案是壓力。早期的輻射壓很大。利用
\begin{equation}
p=w\rho,
\end{equation}
常見的情形是
\begin{equation}
w_{\rm rad}=\frac13,\qquad
w_m=0,\qquad
w_{\rm vac}=-1.
\end{equation}
正壓力會使介質趨於均勻，將擾動推開並使介質平滑。在大尺度上，物質可視為沒有壓力，因此不會提供這種支撐。有壓力時，引力受到抑制；壓力變得可以忽略時，引力便占上風。

暗物質對事件發生的先後順序至關重要。最先坍縮形成的是暗物質暈。接著，普通重子物質落入這些暈中；由於重子物質能夠耗散，會聚集到較靠近中心的地方。暗物質不會以同樣的方式耗散，因此仍分布在整個暈中。收縮大約在透明時期之前開始，時間上約早十倍；因此，微波背景並未直接呈現這一過程的起始。

插一句，交代後來的化學演化：幾乎所有比氦重的元素都是恆星演化的產物，而非早期宇宙學的產物。早期宇宙基本上只產生氫、氦，以及少量幾種輕原子核。

天空中觀測到的不是質量密度本身，而是微波背景的溫度各向異性。密度漲落會對它有所貢獻，速度漲落也會。過密區域往往看起來稍冷一些，因為光子離開較深的引力勢阱時會額外損失能量；此外，氣體運動造成的都卜勒頻移也會疊加其上。完整的數值分析相當複雜。較可靠、也較簡單的結論是：宇宙變得透明時，能量密度的變化約為\(10^5\)分之一。

\begin{classroomqa}
\lecturetimestamp{00:20:29}
\audiencequestion{為什麼引力沒有讓這些團塊更早增長？}
\lecturerresponse{因為早期宇宙由輻射主導，而輻射帶有足以使介質均勻化的正壓力。只有當輻射相較於物質變得不重要時，引力才開始主導團塊的增長。}
\end{classroomqa}

\begin{classroomqa}
\lecturetimestamp{00:24:35}
\audiencequestion{坍縮不需要耗散機制嗎？}
\lecturerresponse{普通重子物質要沉降到中心，確實需要耗散。但暗物質不必經過那種碰撞耗散，也能先形成暗物質暈。此外還有一種與膨脹本身有關的宇宙學阻尼，稍後會談到它的重要性。}
\end{classroomqa}

\begin{classroomqa}
\lecturetimestamp{00:30:34}
\audiencequestion{是否也有速度漲落，讓後來的星系能夠獲得角動量？}
\lecturerresponse{是。原初圖樣包含密度漲落與速度漲落。天空中觀測到的溫度圖樣，是由密度變化與氣體運動共同造成的。}
\end{classroomqa}

\section{暴脹子與平臺}

接下來介紹的機制涉及一個標量場\(\phi\)。它的能量密度由三部分組成：
\begin{equation}
\rho_\phi=\frac12\dot\phi^2+\frac12(\nabla\phi)^2+V(\phi).
\end{equation}
時間變化會有所貢獻，空間變化也會；標量場還可以帶有真正的勢能\(V(\phi)\)。這部分並不奇特，希格斯場就是熟悉的例子。關鍵在於勢能的形狀。

勢能在某個範圍內必須相當平坦，只有微微向下的斜坡；接著到了某個臨界值，勢能便陡然下降。它不會恰好降至零，而是降到非常小的值。以普朗克單位表示，當今的真空能量約為
\begin{equation}
10^{-123}.
\end{equation}
平臺上的能量比這大得多，但仍遠低於普朗克尺度。

\begin{figure}[t]
\centering
\includegraphics[width=0.78\textwidth]{lecture_08_figure_03.png}
\caption{定性示意暴脹子勢能：寬闊的平臺、陡降的斜坡，以及低處的極小值。}
\lectureframe{00:43:23}
\end{figure}

這張圖描述的不是空間中的位置，而是場的取值。說宇宙從平臺高處開始，只表示在一個極大的區域內，場值都位於那個高處。起初\(\phi\)在不同位置有些變化並不重要，因為暴脹會把這些變化大幅拉伸。膨脹得夠久之後，任何局部區域看起來都幾乎均勻。\editorialnote{黑板畫面中的座標軸標籤無法辨認，但口頭討論清楚指出，縱軸是勢能\(V(\phi)\)，橫軸是場值\(\phi\)。}

只要場停留在勢能寬闊且緩緩傾斜的部分，勢能的作用就像真空能。若\(\phi\)完全不變，它的作用便會和真空能完全相同。先不計省略的數值因子，講課中給出的關係就是
\begin{equation}
H\propto \sqrt{V(\phi)},
\end{equation}
由此得到近似指數膨脹定律
\begin{equation}
a(t)\propto e^{Ht}.
\end{equation}
如今能觀測到的幾乎一切，都發生在場越過斜坡邊緣之後。場從邊緣下降時，平臺上大量類真空能量便轉化為輻射、普通物質、暗物質，以及今天觀測到的微量殘餘真空能。

\begin{classroomqa}
\lecturetimestamp{00:44:07}
\audiencequestion{\(\phi\)是什麼場？這裡仍然使用愛因斯坦方程嗎？}
\lecturerresponse{是。\(\phi\)是時空上的場，而且仍然使用同一個弗里德曼方程。在這種情況下，只要場變化得夠慢，方程右側的能量密度就主要由\(V(\phi)\)提供。}
\end{classroomqa}

\begin{classroomqa}
\lecturetimestamp{00:47:04}
\audiencequestion{類真空的暴脹子能量如何轉化為物質和輻射？}
\lecturerresponse{場滾到谷底並開始振盪時，這些振盪就是場的量子。若這些量子不穩定，便可能衰變成光子、電子、正電子和其他熟悉的粒子。}
\end{classroomqa}

\section{慢滾與哈伯摩擦}

場越過斜坡邊緣後，為什麼不會一直來回振盪？為什麼不會來回擺動，甚至再度爬上去？這裡給出的答案是宇宙學阻尼機制。它並非真正的微觀摩擦，但在運動方程中的作用類似摩擦。

忽略空間梯度，將場視為均勻。如此一來，動能項與勢能項看起來就像一個具有\(T-V\)動力學的單一自由度；但因為它們是能量密度，必須乘上物理體積\(a^3\)。因此可得到如下形式的均勻有效拉格朗日量：
\begin{equation}
L_{\mathrm{hom}}=a^3(t)\left[\frac12\dot\phi^2-V(\phi)\right].
\end{equation}
\editorialnote{黑板最上方的一行被部分裁掉。這裡只在足以重現下方清楚寫出的運動方程的範圍內，重建這個均勻拉格朗日量。}

黑板上清楚寫出的方程是
\begin{equation}
\frac{d}{dt}\left(a^3\dot\phi\right)=-a^3\frac{\partial V}{\partial\phi}.
\end{equation}
展開時間導數後得到
\begin{equation}
a^3\ddot\phi+3a^2\dot a\,\dot\phi=-a^3\frac{\partial V}{\partial\phi},
\end{equation}
再除以\(a^3\)，便有
\begin{equation}
\ddot\phi+3H\dot\phi=-\frac{\partial V}{\partial\phi}.
\end{equation}
講者在口頭推導中明確更正了負號：力等於勢能對場的導數取負。

\begin{figure}[t]
\centering
\includegraphics[width=0.78\textwidth]{lecture_08_figure_05.png}
\caption{均勻暴脹子的運動方程，顯示哈伯摩擦項\(3H\dot\phi\)。}
\lectureframe{00:57:10}
\end{figure}

\(3H\dot\phi\)這一項稱為哈伯摩擦，完全來自對因子\(a^3\)微分。不過，就動力學而言，它看起來就像石塊落入水中時受到的阻力項。方程中有力項、加速度項，以及一個與速度成正比的項。場很快便達到終端速度，接著沿著平臺緩慢移動，這就是慢滾階段。斜坡平緩、場移動緩慢時，勢能變化也很慢，因此膨脹維持近似指數形式。

場在極小值附近振盪。這些振盪就是暴脹子量子。若這些量子不穩定，便會發生衰變，平臺上的能量也就藉此轉化為後期宇宙中的熱物質與輻射。

\begin{classroomqa}
\lecturetimestamp{00:50:48}
\audiencequestion{場為什麼不會一直振盪下去？阻尼從何而來？}
\lecturerresponse{膨脹本身會產生類似摩擦的項。均勻拉格朗日量中的因子\(a^3\)會產生\(3H\dot\phi\)項，其作用就像黏滯阻力。}
\end{classroomqa}

\section{量子漲落與越過視界}

現在看來似乎有個問題：若暴脹使宇宙變得極其平滑，為什麼沒有把星系所需的種子也抹掉？答案在於量子力學。留下來的漲落並非不完全平滑所殘留的痕跡，而是重新產生的。

先暫時不管膨脹。即使場幾乎完全均勻，量子力學仍不允許\(\phi\)與\(\dot\phi\)同時具有完全確定的值。基態仍然存在擾動，這和普通諧振子具有零點能是同一個原因。在日常情況下，這些漲落很難觀測，因為它們波長短、振盪快，觀測時會對它們取平均。

膨脹改變了情況。在共動座標中，場模態的座標波長大致不變，物理波長則隨尺度因子增長：
\begin{equation}
\lambda_{\rm phys}=a(t)\,\Delta x.
\end{equation}
暴脹期間，物理哈伯尺度近似固定，
\begin{equation}
d_H\sim H^{-1},
\end{equation}
因此相應的共動視界會以如下方式縮小：
\begin{equation}
d_{H,\mathrm{coord}}\sim \frac{1}{aH}.
\end{equation}

繩子的比喻有助於理解模態如何凍結。波會振盪，是因為繩子的每一小段都受到相鄰部分的牽引。但若波長大到和視界相當，相關的相鄰部分便超出因果接觸範圍，恢復力隨之消失。模態一旦長於視界，就會凍結。

這種情況不斷重演：長波長先凍結，短波長後凍結。即使某個模態起初的波長極短，最終也會被拉伸到越過視界。因此，暴脹並非只會抹去結構；它還會持續產生凍結的場變化。

即使暴脹持續很久，這個過程也不會停止。舊的凍結模態被拉伸到極大尺度，新的模態則不斷凍結。就宇宙學所關心的意義而言，累積而成的圖樣在統計上趨於穩定。近似尺度不變譜便由此而來：同一個越過視界的過程在接續的對數尺度上反覆發生，使相鄰波長帶的振幅大致相同。

場滾下平臺後，視界再度增長。先前位於暫時暴脹視界之外的模態，之後會重新進入視界。通常會說，這些模態在暴脹期間離開視界，之後又重新進入。此時它們已被拉伸至宏觀尺度，因此再次開始振盪時，呈現的是非常大的結構。

\begin{classroomqa}
\lecturetimestamp{01:12:02}
\audiencequestion{在時空圖上，視界看起來是什麼樣子？}
\lecturerresponse{指數膨脹期間，到視界的物理距離大致維持在\(c/H\)；但在共動座標圖上，由於座標尺度不斷被拉伸，這段固定的物理距離看起來會愈來愈短。換成共動距離來說，視界會按\(1/(aH)\)縮小。}
\end{classroomqa}

\begin{classroomqa}
\lecturetimestamp{01:15:59}
\audiencequestion{漲落被拉伸時，振幅會如何變化？會形成穩態嗎？}
\lecturerresponse{個別模態會在拉伸時稀釋，但新的模態不斷凍結。結果是形成統計上穩定、近似尺度不變的圖樣，而不是讓漲落完全消失。}
\end{classroomqa}

\section{從凍結的場變化到密度變化}

場在不同位置產生微小變化後，若不同區域越過斜坡邊緣的時間略有差異，這些變化便會轉化為密度差異。沿斜坡稍微更低處的區域會先到達邊緣。若它的真空能轉化為輻射，輻射便會隨膨脹而稀釋。相鄰區域若仍在平臺上，就仍由真空能主導，不會稀釋。等到第二個區域落下時，第一個區域已有一部分能量因紅移而流失。

這就是為什麼場滾動得愈慢，最後形成的密度反差愈大。若場沿著平臺移動得很慢，場值的微小差異便會使相鄰區域到達斜坡邊緣的時間差拉大。較早落下的區域有更多時間稀釋，等到較晚落下的區域才轉化其能量。斜坡愈平緩，終端速度愈小，最後的密度變化也愈大。平臺愈高，漲落通常也愈大。平臺寬度主要決定暴脹持續的時間。

漲落資料嚴格限制的，是能量高度與斜率的一種組合。若能量尺度夠高，也應會有可觀測的引力波貢獻。這裡給出的數值界限只是粗略估計。大致而言，若平臺上的能量密度遠高於約\(10^{-14}\)個普朗克單位，微波背景中就會出現過多的引力波訊號；若遠低於約\(10^{-16}\)個普朗克單位，實際上就很難觀測到這種訊號。\editorialnote{講者明確表示這些數值只是近似值，並說自己記不清確切數字。此處僅保留其數量級估計。}

單位依照自然單位制的慣例說明。質量為\(m\)的粒子，其康普頓波長為
\begin{equation}
\lambda_C=\frac{\hbar}{mc}.
\end{equation}
若\(c=\hbar=1\)，質量與長度的倒數便具有相同單位。能量密度是單位體積內的能量，因此其單位是質量的四次方：
\begin{equation}
[\rho]=(\mathrm{mass})^4.
\end{equation}

在單一位置觀察時，凍結模態的累積看起來像隨機漫步。空間中不同位置各自經歷不同的隨機漫步，正是這種差異將量子雜訊轉化為空間變化。

\begin{classroomqa}
\lecturetimestamp{01:33:04}
\audiencequestion{\(\phi\)會隨機漫步嗎？這是否使更早的起始狀態變得無關緊要？}
\lecturerresponse{在任一固定位置，相繼凍結的模態以隨機相位相加，因此\(\phi\)確實會隨機漫步。但這項成果本身也帶來了困難：暴脹不斷以後來的漲落取代早先的資訊，使暴脹之前的階段難以重建。}
\end{classroomqa}

\section{暴脹所隱藏的事物}

這個理論在一個特定意義上相當成功：它以同一個底層機制，連結起跨越巨大尺度範圍的漲落。但同一個機制也會沖淡此前的狀態。早期的不規則性被拉伸到難以察覺，之後產生的量子漲落則成為仍然重要的部分。因此，暴脹既引人注目，也令人困擾。它解釋了宇宙的平滑性與漲落譜，卻也抹去了大部分更早的歷史。

這讓人對起始狀態感到不安。問題不在於能量密度非常高，而在於平臺能量看起來既非自然落在普朗克尺度，也非自然地極小。以一般單位來看，它很高；但以普朗克單位來看，它仍然很小。它不像是基礎理論中一個顯而易見的起始位置。

在這個簡單圖像中，也沒有字面上的\(t=0\)。暴脹一旦持續太久，幾乎所有關於更早階段的資訊都會消失。若 e-fold 數只比所需的最小值略高，也許還能留下前一階段的某些痕跡；但一般而言，暴脹會抹去過去。

處理德西特空間時也有同樣的保留。純粹的指數定律適用於暴脹分支的晚期，但精確的德西特解更適合寫成
\begin{equation}
a(t)\propto \cosh(Ht).
\end{equation}
這個解會先收縮到最小尺度，然後再膨脹。收縮分支是否應在宇宙學中按字面理解，仍是個開放問題。

\begin{classroomqa}
\lecturetimestamp{01:37:31}
\audiencequestion{宇宙會不會沒有起點，而指數膨脹一直向過去延續？}
\lecturerresponse{在這個簡單圖像中不會。精確的德西特空間並非在所有時間都永遠指數膨脹；它更像\(\cosh(Ht)\)，包含收縮分支和膨脹分支。較早的那個分支是否應按字面理解，仍是個開放問題。}
\end{classroomqa}

\section{微調、曲率與人擇界限}

粗略來說，重要的參數包括平臺的高度、斜率與寬度\(\Delta\phi\)。斜率會決定終端速度，進而決定漲落大小，因此很重要；寬度則有助於決定暴脹持續多久。一種常見的看法是，基本的場位移不應遠大於普朗克質量，而能量密度不應超過普朗克單位下的一個量級。在這些粗略界限下，能產生足夠暴脹的參數並未遍及整個參數空間。此處引用的估計是，需要微調百分之幾，約為\(2\%\)，才能讓斜率足夠平緩，以產生約\(60\)至\(62\)個 e-fold。

\begin{classroomqa}
\lecturetimestamp{01:46:02}
\audiencequestion{能否用參數空間來理解暴脹？需要多大程度的微調？}
\lecturerresponse{可以。粗略來說，可以考慮平臺的高度、斜率與寬度。講課中提出的估計是，要有足夠的 e-fold 數，需要百分之幾程度的微調，約為\(2\%\)。}
\end{classroomqa}

接著，討論將這種程度不大的微調，與當前宇宙學常數的難題相比；後者的量級約為\(10^{-123}\)。話題轉向人擇問題，因為這兩個難題都可以改寫成：形成結構需要具備什麼條件？

在牛頓圖像中，正的宇宙學常數會造成一種與距離成正比的普遍斥力。若它大得多，不必等到今天把星系撕裂，就足以造成致命影響；只要大到能阻止原初\(10^{-5}\)漲落凝聚成星系即可。講者舉的例子是，若其數值約比觀測值高兩個數量級，就足以造成這種結果。這就是溫伯格論證的要點：宇宙學常數不必恰好為零，只須小到足以讓結構形成。

負曲率也會產生同方向的作用。用先前講課中的牛頓語言來說，\(k<0\)表示退行速度大於逃逸速度。向外推力若過強，也會使微小的原初過密區無法發展成星系。暴脹會稀釋這種曲率。引用的下界約為\(59.5\)個 e-fold。約\(62\)個 e-fold 則被視為一段高於此門檻、具有代表性的暴脹歷史。

這改變了機率問題。若問能產生足夠暴脹的參數空間占全部參數空間的多少，答案可能只有百分之幾。若把條件限定在能形成星系、也有觀察者得以提出這個問題的宇宙中，條件機率就可能大得多。討論最後停在這裡：沒有提出最終解答，而是回到微調難題、對這些難題的人擇詮釋，以及一個認知——重大的爭論仍然在此。